\documentclass[11pt]{article}
\usepackage[margin=1in]{geometry}
\usepackage{xcolor}
\usepackage[breakable,listings]{tcolorbox}
\tcbuselibrary{listings}
\usepackage{hyperref}
\usepackage{enumitem}
\usepackage{listings}
\usepackage{verbatim}
\usepackage{amsmath}
\usepackage{graphicx}
\usepackage{booktabs}
\usepackage{float}

% Define colors
\definecolor{osaorange}{RGB}{220,115,38}
\definecolor{aiblue}{RGB}{30,144,255}
\definecolor{codebackground}{RGB}{248,248,248}

% Configure hyperref
\hypersetup{
    colorlinks=true,
    linkcolor=blue,
    urlcolor=blue,
    citecolor=blue
}

% Configure listings for code display
\lstset{
  basicstyle=\small\ttfamily,
  breaklines=true,
  breakatwhitespace=true,
  postbreak=\mbox{\textcolor{gray}{$\hookrightarrow$}\space},
  columns=flexible,
  keepspaces=true,
  showstringspaces=false
}

% Code environment using tcolorbox with listings
\newtcblisting{rcode}{
  colback=white,
  colframe=gray!50,
  boxrule=0.5pt,
  arc=0pt,
  left=3pt,
  right=3pt,
  top=3pt,
  bottom=3pt,
  width=\textwidth,
  breakable,
  listing only,
  listing options={
    basicstyle=\small\ttfamily,
    breaklines=true,
    breakatwhitespace=true,
    postbreak=\mbox{\textcolor{gray}{$\hookrightarrow$}\space},
    columns=flexible,
    keepspaces=true,
    showstringspaces=false
  }
}

\title{}
\author{}
\date{}

\begin{document}

% Header box
\begin{tcolorbox}[colback=osaorange,colframe=black,boxrule=2pt,arc=0pt,left=3pt,right=3pt,top=6pt,bottom=6pt,boxsep=2pt]
\centering
\textcolor{white}{\textbf{\Large H524 Introduction to Biostatistics}}\\
\textcolor{white}{\textbf{\Large Week 8 Lab: Contingency Tables and Chi-Square Tests}}\\
\textcolor{white}{\textbf{\Large Fall 2025}}
\end{tcolorbox}

\vspace{8pt}

\noindent\textbf{Format:} Online\\
\textbf{Tools Required:} R, RStudio\\
\textbf{Estimated Time:} 90 minutes

\section{Learning Objectives}

By the end of this lab, you will be able to:
\begin{enumerate}
\item Perform chi-square goodness-of-fit tests in R
\item Conduct chi-square tests of independence for contingency tables
\item Use Fisher's exact test for small samples
\item Calculate and interpret odds ratios and relative risks
\item Create and interpret mosaic plots and grouped bar charts
\item Determine when to use different categorical data tests
\item \textcolor{aiblue}{\textbf{Use AI tools to help with categorical data analysis}}
\end{enumerate}

\section{Part 1: Chi-Square Goodness-of-Fit Test}

The chi-square goodness-of-fit test determines whether observed frequencies match expected frequencies under a specified distribution.

\subsection{Exercise 1.1: Blood Type Distribution}

\textbf{Scenario:} The U.S. population has the following blood type distribution: O (45\%), A (40\%), B (11\%), AB (4\%). We sample 200 hospital patients.

\begin{rcode}
# Observed counts from hospital sample
observed <- c(O=85, A=78, B=28, AB=9)
total <- sum(observed)

cat("Total sample size:", total, "\n")
cat("Observed counts:\n")
print(observed)

# Expected proportions (from population)
expected_props <- c(O=0.45, A=0.40, B=0.11, AB=0.04)

# Chi-square goodness-of-fit test
chi_result <- chisq.test(observed, p=expected_props)

cat("\nChi-square test results:\n")
print(chi_result)

# Extract components
cat("\nTest statistic:", round(chi_result$statistic, 3), "\n")
cat("Degrees of freedom:", chi_result$parameter, "\n")
cat("P-value:", round(chi_result$p.value, 4), "\n")

# Expected counts
cat("\nExpected counts:\n")
print(round(chi_result$expected, 2))

# Compare observed vs expected
comparison <- data.frame(
  BloodType = names(observed),
  Observed = as.numeric(observed),
  Expected = round(chi_result$expected, 2),
  Difference = as.numeric(observed) - round(chi_result$expected, 2)
)
print(comparison)
\end{rcode}

\textbf{Interpretation:}
\begin{itemize}
\item If p-value $>$ 0.05: Observed distribution is consistent with expected (fail to reject $H_0$)
\item If p-value $<$ 0.05: Observed distribution differs significantly from expected (reject $H_0$)
\end{itemize}

\subsection{Exercise 1.2: Visualizing Goodness-of-Fit}

\begin{rcode}
# Create visualization
library(graphics)

# Prepare data for plotting
blood_data <- rbind(
  Observed = as.numeric(observed),
  Expected = chi_result$expected
)
colnames(blood_data) <- names(observed)

# Grouped barplot
barplot(blood_data, beside=TRUE,
        col=c("steelblue", "coral"),
        legend=rownames(blood_data),
        main="Blood Type Distribution: Observed vs Expected",
        xlab="Blood Type", ylab="Count",
        args.legend=list(x="topright"))

# Add chi-square test result as text
text(x=6, y=max(blood_data)*0.9,
     labels=paste0("Chi-square = ", round(chi_result$statistic, 2),
                   "\nP-value = ", round(chi_result$p.value, 3)),
     pos=4, cex=0.9)
\end{rcode}

\subsection{Exercise 1.3: Mendelian Genetics Example}

\textbf{Scenario:} Mendel's theory predicts a 3:1 ratio of dominant to recessive traits. We observe 315 dominant and 108 recessive offspring.

\begin{rcode}
# Observed counts
observed <- c(Dominant=315, Recessive=108)
total <- sum(observed)

# Expected proportions (3:1 ratio = 0.75:0.25)
expected_props <- c(Dominant=0.75, Recessive=0.25)

# Chi-square test
chi_result <- chisq.test(observed, p=expected_props)

cat("Mendelian Genetics Test\n")
cat("========================\n")
cat("Observed: Dominant =", observed[1], ", Recessive =", observed[2], "\n")
cat("Expected proportions: 3:1 (75%:25%)\n\n")

cat("Chi-square statistic:", round(chi_result$statistic, 4), "\n")
cat("P-value:", round(chi_result$p.value, 4), "\n\n")

if (chi_result$p.value > 0.05) {
  cat("Conclusion: Data consistent with Mendelian 3:1 ratio\n")
} else {
  cat("Conclusion: Data NOT consistent with Mendelian 3:1 ratio\n")
}
\end{rcode}

\subsection{Exercise 1.4: Practice - Die Fairness}

\textbf{Your turn:} Test if a die is fair. In 120 rolls, we observe: 1(18), 2(22), 3(17), 4(23), 5(19), 6(21).

\begin{rcode}
# Enter the data
observed <- c(one=18, two=22, three=17, four=23, five=19, six=21)

# Expected proportions for fair die (all equal)
expected_props <- rep(1/6, 6)

# Perform chi-square test
chi_result <- chisq.test(observed, p=expected_props)

# Display results
cat("Die Fairness Test\n")
cat("=================\n")
print(chi_result)

# Conclusion
cat("\nConclusion at alpha=0.05:\n")
if (chi_result$p.value > 0.05) {
  cat("The die appears to be fair (p =", round(chi_result$p.value, 3), ")\n")
} else {
  cat("The die appears to be biased (p =", round(chi_result$p.value, 3), ")\n")
}
\end{rcode}

\section{Part 2: Chi-Square Test of Independence}

The chi-square test of independence assesses whether two categorical variables are associated.

\subsection{Exercise 2.1: $2 \times 2$ Contingency Table}

\textbf{Scenario:} Study of smoking and lung cancer. 100 cancer patients, 100 controls.

\begin{rcode}
# Create contingency table
data <- matrix(c(85, 55, 15, 45), nrow=2, byrow=TRUE)
rownames(data) <- c("Smoker", "Non-smoker")
colnames(data) <- c("Cancer", "No Cancer")

cat("Contingency Table:\n")
print(data)
cat("\nTable with margins:\n")
print(addmargins(data))

# Chi-square test
chi_result <- chisq.test(data)

cat("\nChi-square Test Results:\n")
cat("========================\n")
cat("Test statistic:", round(chi_result$statistic, 3), "\n")
cat("Degrees of freedom:", chi_result$parameter, "\n")
cat("P-value:", format.pval(chi_result$p.value, digits=3), "\n")

# Expected counts
cat("\nExpected counts (under independence):\n")
print(round(chi_result$expected, 2))

# Check assumption
all_expected_ge_5 <- all(chi_result$expected >= 5)
cat("\nAll expected counts >= 5?", all_expected_ge_5, "\n")
if (!all_expected_ge_5) {
  cat("WARNING: Chi-square assumptions may be violated!\n")
  cat("Consider Fisher's exact test instead.\n")
}
\end{rcode}

\subsection{Exercise 2.2: Calculating Effect Sizes}

\begin{rcode}
# Using the smoking/cancer data from above

# Calculate Relative Risk
p_cancer_smoker <- data[1,1] / sum(data[1,])
p_cancer_nonsmoker <- data[2,1] / sum(data[2,])

relative_risk <- p_cancer_smoker / p_cancer_nonsmoker

cat("Risk Calculations:\n")
cat("==================\n")
cat("Risk of cancer among smokers:", round(p_cancer_smoker, 4), "\n")
cat("Risk of cancer among non-smokers:", round(p_cancer_nonsmoker, 4), "\n")
cat("Relative Risk:", round(relative_risk, 3), "\n\n")

# Calculate Odds Ratio
odds_ratio <- (data[1,1] * data[2,2]) / (data[1,2] * data[2,1])
cat("Odds Ratio:", round(odds_ratio, 3), "\n\n")

# Confidence interval for OR (using log transformation)
a <- data[1,1]
b <- data[1,2]
c <- data[2,1]
d <- data[2,2]

log_or <- log(odds_ratio)
se_log_or <- sqrt(1/a + 1/b + 1/c + 1/d)
ci_log <- log_or + c(-1.96, 1.96) * se_log_or
ci_or <- exp(ci_log)

cat("95% Confidence Interval for OR:\n")
cat("(", round(ci_or[1], 3), ", ", round(ci_or[2], 3), ")\n\n", sep="")

# Interpretation
cat("Interpretation:\n")
cat("- Smokers are", round(relative_risk, 2),
    "times more likely to have lung cancer\n")
cat("- The odds of cancer among smokers are", round(odds_ratio, 2),
    "times the odds among non-smokers\n")
cat("- We're 95% confident the true OR is between",
    round(ci_or[1], 2), "and", round(ci_or[2], 2), "\n")
\end{rcode}

\subsection{Exercise 2.3: Visualizing $2 \times 2$ Tables}

\begin{rcode}
# Multiple visualization approaches

par(mfrow=c(2,2))  # 2x2 plot layout

# 1. Mosaic plot
mosaicplot(data, main="Mosaic Plot",
           color=c("red", "lightblue"),
           xlab="Smoking Status", ylab="Cancer Status")

# 2. Grouped barplot (counts)
barplot(data, beside=TRUE, legend=TRUE,
        col=c("darkred", "darkgreen"),
        main="Counts by Group",
        xlab="Cancer Status", ylab="Count",
        args.legend=list(x="topright", title="Smoking"))

# 3. Row proportions (shows risk by smoking status)
prop_row <- prop.table(data, margin=1)
barplot(prop_row, beside=TRUE, legend=TRUE,
        col=c("darkred", "darkgreen"),
        main="Cancer Rate by Smoking Status",
        xlab="Cancer Status", ylab="Proportion",
        args.legend=list(x="topright", title="Smoking"))

# 4. Column proportions (shows smoking distribution by cancer status)
prop_col <- prop.table(data, margin=2)
barplot(prop_col, col=c("darkred", "darkgreen"),
        main="Smoking Distribution by Cancer Status",
        xlab="Cancer Status", ylab="Proportion",
        legend=TRUE,
        args.legend=list(x="topleft", title="Smoking"))

par(mfrow=c(1,1))  # Reset layout
\end{rcode}

\subsection{Exercise 2.4: Larger Contingency Tables}

\textbf{Scenario:} $3 \times 3$ table of age group vs disease severity.

\begin{rcode}
# Create 3x3 table
data <- matrix(c(45, 20, 5,   # Young
                 30, 35, 15,  # Middle
                 15, 25, 10), # Older
               nrow=3, byrow=TRUE)
rownames(data) <- c("Young (<40)", "Middle (40-60)", "Older (>60)")
colnames(data) <- c("Mild", "Moderate", "Severe")

cat("Disease Severity by Age Group:\n")
print(data)
cat("\nWith margins:\n")
print(addmargins(data))

# Chi-square test
chi_result <- chisq.test(data)

cat("\nChi-square Test:\n")
cat("================\n")
cat("Test statistic:", round(chi_result$statistic, 3), "\n")
cat("Degrees of freedom:", chi_result$parameter, "\n")
cat("P-value:", format.pval(chi_result$p.value, digits=4), "\n\n")

# Expected counts
cat("Expected counts:\n")
print(round(chi_result$expected, 2))

# Standardized residuals (show which cells contribute most)
std_residuals <- chi_result$residuals
cat("\nStandardized Residuals:\n")
cat("(Values > 2 or < -2 indicate major contributors to chi-square)\n")
print(round(std_residuals, 2))

# Visualization
mosaicplot(data, main="Disease Severity by Age Group",
           color=c("lightgreen", "yellow", "red"),
           xlab="Age Group", ylab="Severity")
\end{rcode}

\textbf{Note on interpretation:} For larger tables, significant chi-square tells us there's an association, but doesn't specify the pattern. Examine standardized residuals to identify which cells differ most from expected.

\subsection{Exercise 2.5: Practice - Treatment Efficacy}

\textbf{Your turn:} Clinical trial with new drug vs placebo.

\begin{rcode}
# Enter data
data <- matrix(c(48, 12,   # Drug: cured, not cured
                 32, 28),  # Placebo: cured, not cured
               nrow=2, byrow=TRUE)
rownames(data) <- c("Drug", "Placebo")
colnames(data) <- c("Cured", "Not Cured")

# Display table
cat("Treatment Efficacy Study:\n")
print(addmargins(data))

# Chi-square test
chi_result <- chisq.test(data)
cat("\nChi-square test:\n")
print(chi_result)

# Calculate effect sizes
rr <- (data[1,1]/sum(data[1,])) / (data[2,1]/sum(data[2,]))
or <- (data[1,1]*data[2,2]) / (data[1,2]*data[2,1])

cat("\nEffect Sizes:\n")
cat("Relative Risk:", round(rr, 3), "\n")
cat("Odds Ratio:", round(or, 3), "\n")

# Visualize
barplot(prop.table(data, margin=1), beside=TRUE,
        col=c("steelblue", "coral"),
        legend=TRUE, main="Cure Rates by Treatment",
        xlab="Outcome", ylab="Proportion")
\end{rcode}

\section{Part 3: Fisher's Exact Test}

Fisher's exact test is used when chi-square assumptions are violated (small expected counts).

\subsection{Exercise 3.1: Small Sample Example}

\textbf{Scenario:} Rare disease study with small sample size.

\begin{rcode}
# Small sample data
data <- matrix(c(8, 2,   # Exposed: disease, no disease
                 3, 7),  # Not exposed: disease, no disease
               nrow=2, byrow=TRUE)
rownames(data) <- c("Exposed", "Not Exposed")
colnames(data) <- c("Disease", "No Disease")

cat("Small Sample Study:\n")
print(data)
print(addmargins(data))

# Check chi-square assumptions
chi_result <- chisq.test(data)
cat("\nExpected counts:\n")
print(round(chi_result$expected, 2))

cat("\nNote: Some expected counts < 5, so Fisher's exact test preferred\n\n")

# Fisher's exact test
fisher_result <- fisher.test(data)

cat("Fisher's Exact Test Results:\n")
cat("=============================\n")
print(fisher_result)

cat("\nKey results:\n")
cat("P-value (two-sided):", round(fisher_result$p.value, 4), "\n")
cat("Odds Ratio estimate:", round(fisher_result$estimate, 3), "\n")
cat("95% CI for OR: (",
    round(fisher_result$conf.int[1], 3), ", ",
    round(fisher_result$conf.int[2], 3), ")\n", sep="")
\end{rcode}

\subsection{Exercise 3.2: Comparing Fisher's vs Chi-Square}

\begin{rcode}
# For comparison purposes with same data

cat("Comparison of Tests:\n")
cat("====================\n\n")

# Chi-square test (for comparison, though not ideal here)
chi_result <- chisq.test(data)
cat("Chi-square test:\n")
cat("  Test statistic:", round(chi_result$statistic, 3), "\n")
cat("  P-value:", round(chi_result$p.value, 4), "\n\n")

# Fisher's exact test
fisher_result <- fisher.test(data)
cat("Fisher's exact test:\n")
cat("  P-value:", round(fisher_result$p.value, 4), "\n")
cat("  Odds Ratio:", round(fisher_result$estimate, 3), "\n")
cat("  95% CI:", round(fisher_result$conf.int, 3), "\n\n")

cat("Note: Fisher's exact test is more reliable for small samples\n")
cat("Chi-square may give misleading results when expected counts < 5\n")
\end{rcode}

\subsection{Exercise 3.3: Practice - Medical Device Study}

\begin{rcode}
# Very small pilot study
data <- matrix(c(5, 1,   # New device: success, failure
                 2, 6),  # Standard: success, failure
               nrow=2, byrow=TRUE)
rownames(data) <- c("New Device", "Standard")
colnames(data) <- c("Success", "Failure")

# Use Fisher's exact test
fisher_result <- fisher.test(data)

cat("Medical Device Pilot Study:\n")
print(data)
cat("\nFisher's Exact Test:\n")
print(fisher_result)

# Interpretation
if (fisher_result$p.value < 0.05) {
  cat("\nConclusion: Significant difference between devices (p < 0.05)\n")
} else {
  cat("\nConclusion: No significant difference found (p >= 0.05)\n")
}
\end{rcode}

\section{Part 4: McNemar's Test for Paired Data}

McNemar's test is used when observations are paired or matched.

\subsection{Exercise 4.1: Before-After Treatment Study}

\textbf{Scenario:} 50 patients tested before and after treatment for infection.

\begin{rcode}
# Create paired data table
# Rows = Before, Columns = After
data <- matrix(c(10, 20,   # Before Positive: After Pos, After Neg
                 5, 15),   # Before Negative: After Pos, After Neg
               nrow=2, byrow=TRUE)
rownames(data) <- c("Before: Positive", "Before: Negative")
colnames(data) <- c("After: Positive", "After: Negative")

cat("Before-After Study (Paired Data):\n")
print(data)
cat("\nTotal patients:", sum(data), "\n")

# Count discordant pairs
cat("\nDiscordant pairs:\n")
cat("  Positive -> Negative:", data[1,2], "\n")
cat("  Negative -> Positive:", data[2,1], "\n")

# McNemar's test
mcnemar_result <- mcnemar.test(data)

cat("\nMcNemar's Test Results:\n")
cat("=======================\n")
print(mcnemar_result)

# Interpretation
cat("\nInterpretation:\n")
if (mcnemar_result$p.value < 0.05) {
  cat("There IS a significant change in infection status after treatment\n")
  if (data[1,2] > data[2,1]) {
    cat("Treatment appears effective (more patients became negative)\n")
  }
} else {
  cat("No significant change in infection status after treatment\n")
}
\end{rcode}

\subsection{Exercise 4.2: Matched Case-Control Study}

\begin{rcode}
# Matched pairs: cases and controls
# Each case matched with a control on age, sex, etc.
data <- matrix(c(15, 8,    # Case exposed: Control exposed, not exposed
                 5, 30),   # Case not exposed: Control exposed, not exposed
               nrow=2, byrow=TRUE)
rownames(data) <- c("Case: Exposed", "Case: Not Exposed")
colnames(data) <- c("Control: Exposed", "Control: Not Exposed")

cat("Matched Case-Control Study:\n")
print(data)

# NOTE: Standard chi-square would be WRONG here (violates independence)
cat("\nWARNING: Cannot use standard chi-square test!\n")
cat("Pairs are matched (not independent)\n\n")

# Correct test: McNemar's
mcnemar_result <- mcnemar.test(data)

cat("McNemar's Test (correct for matched data):\n")
print(mcnemar_result)

# Calculate Odds Ratio for matched pairs
or_matched <- data[1,2] / data[2,1]
cat("\nOdds Ratio (matched pairs):", round(or_matched, 3), "\n")
\end{rcode}

\section{Part 5: Comprehensive Example}

Let's work through a complete analysis with multiple approaches.

\subsection{Exercise 5.1: Vaccine Efficacy Study}

\textbf{Scenario:} Randomized trial of new vaccine. 500 vaccinated, 500 placebo. Follow for disease development.

\begin{rcode}
# Create data
data <- matrix(c(30, 470,    # Vaccinated: disease, no disease
                 85, 415),   # Placebo: disease, no disease
               nrow=2, byrow=TRUE)
rownames(data) <- c("Vaccinated", "Placebo")
colnames(data) <- c("Disease", "No Disease")

cat("========================================\n")
cat("Vaccine Efficacy Study Analysis\n")
cat("========================================\n\n")

# Display data
cat("Contingency Table:\n")
print(data)
print(addmargins(data))

# 1. Chi-square test
cat("\n\n1. CHI-SQUARE TEST OF INDEPENDENCE\n")
cat("------------------------------------\n")
chi_result <- chisq.test(data)
print(chi_result)

cat("\nExpected counts:\n")
print(round(chi_result$expected, 2))

# 2. Effect sizes
cat("\n\n2. EFFECT SIZES\n")
cat("----------------\n")

# Incidence rates
inc_vax <- data[1,1] / sum(data[1,])
inc_placebo <- data[2,1] / sum(data[2,])

cat("Disease incidence:\n")
cat("  Vaccinated:", round(inc_vax * 100, 2), "%\n")
cat("  Placebo:", round(inc_placebo * 100, 2), "%\n\n")

# Relative Risk
rr <- inc_vax / inc_placebo
cat("Relative Risk:", round(rr, 3), "\n")

# Vaccine Efficacy
ve <- (1 - rr) * 100
cat("Vaccine Efficacy:", round(ve, 2), "%\n\n")

# Odds Ratio
or <- (data[1,1] * data[2,2]) / (data[1,2] * data[2,1])
cat("Odds Ratio:", round(or, 3), "\n")

# 95% CI for OR
a <- data[1,1]; b <- data[1,2]; c <- data[2,1]; d <- data[2,2]
log_or <- log(or)
se_log_or <- sqrt(1/a + 1/b + 1/c + 1/d)
ci_or <- exp(log_or + c(-1.96, 1.96) * se_log_or)
cat("95% CI for OR: (", round(ci_or[1], 3), ", ",
    round(ci_or[2], 3), ")\n\n", sep="")

# 3. Interpretation
cat("\n3. INTERPRETATION\n")
cat("------------------\n")
cat("Statistical significance: p", ifelse(chi_result$p.value < 0.001,
    "< 0.001", paste("=", round(chi_result$p.value, 4))), "\n")
cat("  -> Strong evidence of association between vaccination and disease\n\n")

cat("Vaccine efficacy:", round(ve, 1), "%\n")
cat("  -> Vaccine reduces disease risk by", round(ve, 1), "%\n\n")

cat("Clinical interpretation:\n")
cat("  The vaccine significantly reduces disease incidence.\n")
cat("  Vaccinated individuals have", round((1-rr)*100, 1),
    "% lower risk than unvaccinated.\n")

# 4. Visualization
cat("\n\n4. VISUALIZATION\n")
cat("-----------------\n")

par(mfrow=c(1,2))

# Incidence comparison
inc_data <- c(inc_vax, inc_placebo) * 100
names(inc_data) <- c("Vaccinated", "Placebo")
barplot(inc_data, col=c("steelblue", "coral"),
        main="Disease Incidence by Group",
        ylab="Incidence (%)", ylim=c(0, max(inc_data)*1.2))
text(x=c(0.7, 1.9), y=inc_data + 1, labels=paste0(round(inc_data, 1), "%"))

# Mosaic plot
mosaicplot(data, main="Disease Status by Vaccination",
           color=c("red", "lightgreen"),
           xlab="Vaccination Status", ylab="Disease Status")

par(mfrow=c(1,1))

cat("\nAnalysis complete!\n")
\end{rcode}

\section{Practice Problems}

\subsection{Problem 1: Educational Attainment}

Test if the distribution of highest degree earned in a sample matches national proportions: No HS (12\%), HS (30\%), Bachelor's (35\%), Graduate (23\%). Sample: No HS (8), HS (42), Bachelor's (55), Graduate (35).

\begin{rcode}
# YOUR CODE HERE
# Hint: Use chisq.test() with observed counts and expected proportions
\end{rcode}

\vspace{2cm}

\subsection{Problem 2: Drug Side Effects}

$2 \times 2$ table: Drug A vs Drug B, Side effects vs No side effects. Drug A: 15 with side effects, 85 without. Drug B: 28 with side effects, 72 without.

\begin{enumerate}[label=\alph*)]
\item Perform chi-square test
\item Calculate and interpret relative risk
\item Calculate and interpret odds ratio with 95\% CI
\item Create a visualization
\end{enumerate}

\vspace{2cm}

\subsection{Problem 3: Coffee and Heart Disease}

Small pilot study: Coffee drinkers vs non-drinkers, Heart disease vs none. Coffee (disease=4, none=8), No coffee (disease=1, none=7).

\begin{enumerate}[label=\alph*)]
\item Check if chi-square assumptions are met
\item Perform Fisher's exact test
\item Interpret the results
\end{enumerate}

\vspace{2cm}

\subsection{Problem 4: Depression Treatment}

Paired study: 40 patients assessed for depression before and after therapy. Before depressed/After depressed (12), Before depressed/After not (18), Before not/After depressed (3), Before not/After not (7).

\begin{enumerate}[label=\alph*)]
\item Why is McNemar's test appropriate?
\item Perform the test
\item Interpret the results
\end{enumerate}

\vspace{2cm}

\section{Saving Your Work}

\textbf{IMPORTANT:} Save your R script file!

\begin{rcode}
# At the top of your R script, add:
# Name: Your Name
# Date: Today's Date
# Lab: Week 8 - Contingency Tables and Chi-Square Tests

# Save your workspace
save.image("week8_lab.RData")

# To load it later:
# load("week8_lab.RData")

# Save specific objects
# save(data, chi_result, fisher_result, file="week8_results.RData")
\end{rcode}

\section{\textcolor{aiblue}{AI-Enhanced Learning Tips}}

\begin{tcolorbox}[colback=aiblue!10,colframe=aiblue,title=Using AI Tools for Categorical Data]

\textbf{Good prompts for AI assistance:}
\begin{itemize}
\item ``Show me how to calculate odds ratio with 95\% CI in R for a 2x2 table''
\item ``Explain when I should use Fisher's exact test instead of chi-square''
\item ``Create R code for a mosaic plot with custom colors and labels''
\item ``How do I interpret standardized residuals in a 3x3 contingency table?''
\item ``Calculate vaccine efficacy from this contingency table data''
\end{itemize}

\textbf{AI can help with:}
\begin{itemize}
\item Choosing the appropriate test for your data structure
\item Checking assumptions (expected counts, independence)
\item Calculating effect sizes and confidence intervals
\item Creating publication-quality visualizations
\item Interpreting results in context
\item Debugging R code errors
\end{itemize}

\textbf{Remember:}
\begin{itemize}
\item Always verify chi-square assumptions (expected counts $\geq$ 5)
\item Double-check calculations manually for small tables
\item Distinguish between RR (cohort studies) and OR (case-control)
\item Report both p-values and effect sizes
\item Don't confuse association with causation
\end{itemize}
\end{tcolorbox}

\section{Additional Resources}

\begin{itemize}
\item R Documentation:
  \begin{itemize}
  \item \texttt{?chisq.test} -- Chi-square tests
  \item \texttt{?fisher.test} -- Fisher's exact test
  \item \texttt{?mcnemar.test} -- McNemar's test
  \item \texttt{?mosaicplot} -- Mosaic plots for contingency tables
  \end{itemize}
\item Textbook: Pagano \& Gauvreau, Chapters 15-16
\item Online calculator for effect sizes: \url{https://www.medcalc.org/calc/odds_ratio.php}
\item CRAN Task View: Clinical Trial Design: \url{https://cran.r-project.org/web/views/ClinicalTrials.html}
\end{itemize}

\section{Key Takeaways}

\begin{enumerate}
\item \textbf{Chi-square goodness-of-fit:} Tests if one categorical variable follows expected distribution
\item \textbf{Chi-square independence:} Tests association between two categorical variables
\item \textbf{Assumptions matter:} Check expected counts $\geq$ 5; use Fisher's if violated
\item \textbf{Effect sizes essential:} Report RR or OR with confidence intervals, not just p-values
\item \textbf{Paired data:} Use McNemar's test, not standard chi-square
\item \textbf{Interpretation:} Association does not imply causation; consider confounding
\item \textbf{R Implementation:}
  \begin{itemize}
  \item \texttt{chisq.test()} for goodness-of-fit and independence
  \item \texttt{fisher.test()} for small samples
  \item \texttt{mcnemar.test()} for paired data
  \item Manual calculations for RR and OR with CIs
  \end{itemize}
\end{enumerate}

\end{document}
